\section{Formal categories and de Rham complexes (Part L)}\label{sec:categoryL}

\subsection{The missing equation}
We work over a commutative $\mathbf Q$-algebra $A$.  The constructions
also make sense internally in a ringed topos of $\mathbf Q$-algebras,
using sheaf modules and sheaf tensor products.  For formal objects we
retain the system of nilpotent jets as an adic pro-algebra, as the
``pro-anneau'' on scan 162-5/022 and the opened theorem below do.
Passing to an ordinary inverse-limit sheaf, and recovering its prescribed
jets as ordinary ideal quotients, is an additional effectivity assertion
and is not implicit in this convention.  No finite-generation hypothesis
is imposed on the module of one-forms unless we explicitly say so.

\begin{proposition}[Proved: the corrected generators-and-relations statement]\label{cat:L1}
Let $P$ be an $A$-module.  The data defining a degree-one graded derivation
$d$ of $\Lambda_A^\bullet P$ are an additive derivation
$d_0:A\to P$ and an additive map $d_1:P\to\Lambda_A^2P$ satisfying
\[
 d_1(a\omega)=d_0(a)\wedge\omega+a d_1(\omega).
\]
There is a unique such extension $d$.  It is a differential if and only if
\[
 d_1d_0=0\quad\hbox{and}\quad d_2d_1=0,
 \qquad d_2=d|_{\Lambda_A^2P}.
\]
The first equation alone is insufficient, even for a free module of rank
three over $\mathbf Q$.
\end{proposition}
\begin{proof}
The displayed identity is exactly the compatibility with the relations
$a\omega=(a)\cdot\omega$ in the tensor algebra over $A$.  The graded
Leibniz rule also preserves the exterior relations: the derivative of
$\omega\wedge\omega$ is
$d_1\omega\wedge\omega-\omega\wedge d_1\omega=0$ because degrees
$2$ and $1$ commute.  Additivity handles their polarizations.  Thus the
rule descends from the free graded algebra, and it is unique because $A$
and $P$ generate that algebra.  Applying the rule twice, the two middle
terms cancel, giving
$d^2(uv)=d^2(u)v+u d^2(v)$ for homogeneous $u,v$.
Consequently $d^2$ vanishes everywhere exactly when it vanishes on $A$
and on $P$, which are the two stated equations.

For the counterexample take $A=\mathbf Q$,
$P=\mathbf Qx\oplus\mathbf Qy\oplus\mathbf Qz$, $d_0=0$, and
\[
 dx=y\wedge z,\qquad dy=x\wedge y,\qquad dz=0.
\]
Then $d_1d_0=0$, whereas
$d^2x=(x\wedge y)\wedge z=x\wedge y\wedge z\ne0$.
The program \texttt{code/category\_counterexamples.py} checks the
Leibniz identity on all $64$ ordered pairs of exterior basis elements,
as well as this nonzero square, using exact rational arithmetic.
\end{proof}

We call $(P,d)$ satisfying both equations a \emph{Lie coalgebroid}.
This terminology avoids dualizing a flat module that need not be reflexive.
For example, over $A=\mathbf Q[x]$ the flat nonzero module $P=A[1/x]$
has $\operatorname{Hom}_A(P,A)=0$: the image of $1$ under any such map
belongs to $\bigcap_{n\geq0}x^nA=0$, and torsion-freeness then forces the
entire map to vanish.  The complex $(\Lambda P,0)$ is nevertheless
nonzero.  Thus passage to the ordinary module dual cannot prove the
flat case of the theorem.

\begin{proposition}[Proved: the finite-projective Lie-algebroid dictionary]\label{cat:L2}
If $P$ is finite projective and $T=P^\vee$, the differentials of
Proposition~\ref{cat:L1} correspond bijectively to $A$-Lie algebroid
structures on $T$.  Explicitly the anchor
$\theta:T\to\operatorname{Der}_{\mathbf Q}(A,A)$ and bracket are
\begin{align*}
 \theta(X)(a)&=\langle d_0a,X\rangle,\\
 \langle\omega,[X,Y]\rangle
 &=\theta(X)\langle\omega,Y\rangle
   -\theta(Y)\langle\omega,X\rangle
   -\langle d_1\omega,X\wedge Y\rangle.
\end{align*}
Here a Lie algebroid means a $\mathbf Q$-Lie bracket, an $A$-linear
Lie-homomorphism anchor, and
$[X,aY]=a[X,Y]+\theta(X)(a)Y$.
\end{proposition}
\begin{proof}
The twisted linearity of $d_1$ makes the right side of the bracket
formula $A$-linear in $\omega$: the extra terms involving $d_0a$ cancel.
Finite projectivity identifies this functional with a unique element of
$T$.  Substitution gives skew symmetry and the displayed Leibniz
identity.  On a function the equation $d^2a=0$, evaluated at $X,Y$, is
\[
 [\theta(X),\theta(Y)](a)-\theta([X,Y])(a)=0.
\]
Thus the anchor preserves brackets.  Expanding $d^2\omega$ at $X,Y,Z$
and using this anchor identity cancels the terms with two anchors and
those with one anchor; the remaining expression is
\[
 \langle\omega,
 [[X,Y],Z]+[[Y,Z],X]+[[Z,X],Y]\rangle.
\]
The perfect pairing between $P$ and $T$ therefore gives Jacobi.
Conversely, from a Lie algebroid we define $d_0,d_1$ by the two formulas.
They have the required scalar compatibility.  The same two computations,
now in the reverse direction, give $d^2=0$ on $A$ and $P$; the perfect
pairings with $\Lambda^2T$ and $\Lambda^3T$ detect these equations.
Proposition~\ref{cat:L1} gives $d^2=0$ everywhere.  These constructions
are inverse because their values on the algebra generators agree.
\end{proof}

\subsection{Adic formal categories: symmetry and the associated graded}
An affine adic formal category with object algebra $A$ consists of a
commutative complete separated algebra $B$, source and target maps
$s,t:A\to B$, an identity map $e:B\to A$ with $es=et=1$, and a
composition map
\[
 \Delta:B\longrightarrow B\widehat\otimes_{A,t,s} B
\]
satisfying the category identities.  Write $I=\ker e$ and require
$B\simeq\varprojlim_n B/I^{n+1}$.  The tensor product has the total
ideal filtration.  We write $P=I/I^2$; its two $A$-module structures
agree.  This is the ordinary complete-ring formulation when the jet
system is effective in that form.  In the general pro-algebra
formulation, $B$ denotes its compatible system of augmented algebras
$B_n$ with augmentation ideals $I_n$, $I_n^{n+1}=0$, and
$B_m=B_n/I_n^{m+1}$ for $m\leq n$.  The notation $B/I^{n+1}$ means
the prescribed jet $B_n$; composition and completed tensor products
are defined at every total-order jet.  The same category identities
hold there.  In particular $P=\ker(B_1\to A)$ and all assertions about
associated graded modules are assertions about finite jets.  This
formulation includes arbitrary flat $P$ and sheaves of such pro-algebras.

\begin{lemma}[Proved: flat conormal implies the symmetric graded algebra]\label{cat:L3}
For an adic formal category over $\mathbf Q$ with $P$ flat over $A$,
the canonical surjection
\[
 \operatorname{Sym}_A^nP\longrightarrow I^n/I^{n+1}
\]
is an isomorphism for every $n\geq0$.  Each $B/I^{n+1}$ is flat over
$A$ for both $s$ and $t$.
\end{lemma}
\begin{proof}
We prove injectivity inductively.  Degrees zero and one are immediate.
Suppose the result has been proved below $n$, with $n\geq2$.  Reduce
composition to
\[
 R=(B/I^2)\otimes_{A,t,s}(B/I^n).
\]
Let $I_1,I_2$ be the ideals coming from the first and second factors,
and put $J=I_1+I_2$.  Then $I_1^2=0$, $I_2^n=0$, so
$J^{n+1}=0$ and $J^n=I_1I_2^{n-1}$.
The augmentation of the first factor splits as a right $A$-module map.
Its kernel is $P$, and flatness of $P$ makes
$P\otimes_A(I^{n-1}/I^n)\to P\otimes_A(B/I^n)$ injective.
Consequently multiplication identifies
\[
 J^n=P\otimes_A(I^{n-1}/I^n)
     =P\otimes_A\operatorname{Sym}_A^{n-1}P.
\]
The unit identities for composition say that the linear term of
$\Delta(x)$, for $x\in I$, is $x\otimes1+1\otimes x$ modulo $J^2$.
It follows that the composite on products of $n$ elements of $P$ is
\[
 \partial_n:\operatorname{Sym}_A^nP\longrightarrow
       P\otimes_A\operatorname{Sym}_A^{n-1}P,
 \qquad
 \prod_{i=1}^n w_i\longmapsto
       \sum_{i=1}^n w_i\otimes\prod_{j\ne i}w_j.
\]
Indeed, terms with two factors from $I_1$ or $n$ factors from $I_2$
vanish, and all quadratic corrections have total degree at least $n+1$.
If an element maps to zero in $I^n/I^{n+1}$, its composition image is
zero in $J^n$, hence it lies in $\ker\partial_n$.
Multiplication after $\partial_n$ is $n$ times the identity.
Since $n$ is invertible, this kernel is zero.
This proves the asserted isomorphisms.

Every symmetric power of a flat module over a $\mathbf Q$-algebra is
flat: it is the image of the averaging idempotent on the flat tensor
power.  The exact sequences with successive quotients
$I^j/I^{j+1}=\operatorname{Sym}_A^jP$ therefore show by induction that
all the finite quotients of $B$ are flat.  Source and target give the
same structure on each successive quotient, so the argument applies
to both maps.
\end{proof}

\begin{lemma}[Proved: the inverse exists]\label{cat:L4}
Every formal category in Lemma~\ref{cat:L3} is a formal groupoid.
Its inversion is uniquely determined by composition and identity.
\end{lemma}
\begin{proof}
Consider the shear map sending a composable pair of arrows $(g,h)$ to
$(g,hg)$, a pair with common source.  The two formal schemes have
complete total-ideal filtrations.  The flatness conclusions of
Lemma~\ref{cat:L3} identify their associated graded algebras with
$\operatorname{Sym}_A(P\oplus P)$: for finite quotients, tensor the
successive short exact filtration sequences, which remain exact
because their quotients are flat, and then take the inverse limit.
On degree one the shear has the matrix
$\left(\begin{smallmatrix}1&0\\1&1\end{smallmatrix}\right)$,
up to the choice of order of the two factors.  It is invertible.
Since the associated graded algebra is symmetric, the shear is an
isomorphism on the whole associated graded.  A map between complete
separated filtered rings with this property is an isomorphism: induction
on the filtration length, using the successive short exact sequences,
proves it on every quotient, and inverse limits give its inverse.

Apply the inverse shear to $(g,1_{s(g)})$.  It produces functorially an
arrow $h$ with $hg=1_{s(g)}$.  Apply the same construction to $h$ to
get $k$ with $kh=1_{s(h)}$.  Associativity gives
$k=k(hg)=(kh)g=g$, so $gh=1_{s(h)}$ as well.
The resulting two-sided inverse is unique.  Functoriality of the shear
construction makes it a morphism of formal schemes, proving the claim.
\end{proof}

\begin{theorem}[Known, with the exact hypotheses derived here: corrected Quillen theorem]\label{cat:L5}
Over a ringed topos of $\mathbf Q$-algebras, the functor taking an adic
formal category in the preceding pro-algebra sense, with flat $I/I^2$,
to its invariant de Rham complex is
an equivalence with the category of pairs
\[
 (P,d),\qquad P\text{ flat},\quad
 d:\Lambda_A^qP\longrightarrow\Lambda_A^{q+1}P,
 \quad d^2=0,
\]
where $d$ is a degree-one derivation.  The arrow convention on the
formal side is a continuous coordinate-proalgebra map $B\to B'$
over $A$, respecting source, target, identity and composition.
On the differential side it is a map
$\Lambda_A^\bullet P\to\Lambda_A^\bullet P'$ over $A$ commuting
with the differential.  Thus if arrows on the formal side are instead
written as geometric maps $G\to H$, the equivalence is contravariant:
their pullbacks go from the coordinate algebra and invariant forms of
$H$ to those of $G$.  The construction is also functorial for compatible
changes of $A$.  The inverse integrates the
Lie coalgebroid.  Finite projectivity is unnecessary.
\end{theorem}
\begin{proof}[Derivation from the opened theorem]
Beilinson--Drinfeld, \emph{Chiral Algebras},
\S1.4.14(ii) and \S1.4.15, Proposition
\cite{catBeilinsonDrinfeld2004}, identify Lie coalgebroids with the
stated differential graded algebras and prove their equivalence with
formal groupoids having flat conormal and symmetric associated graded.
Lemmas~\ref{cat:L3} and~\ref{cat:L4} verify exactly these extra
properties for the formal categories at issue.  Conversely their
integrated groupoids give the compatible augmentation-ideal jet system
and have conormal $P$, so they belong to our category.  Their construction uses
tensor products, flat resolutions and completion, and is given in a
unital tensor category; it therefore applies internally to modules in
the ringed topos and is functorial in the object algebra.  The invariant
forms restricted to the identity are $\Lambda_A^\bullet P$, so their
functor is the de Rham functor in the statement.  This proves the
claimed form without replacing flatness by finite projectivity.
All identifications of jets in the source are made at finite order;
we do not replace this pro-algebra theorem by an unproved assertion
that ordinary inverse limits of sheaves preserve its quotients.
\end{proof}

The missing equation in Proposition~\ref{cat:L1} is essential to this
theorem.  A degree-one derivation of nonzero square is not a Lie
coalgebroid and cannot be the invariant de Rham complex of a formal
category.

\subsection{Coefficients and the crystalline starting theorem}
For a Lie coalgebroid $(P,d)$, a connection on an $A$-module $M$ is an
additive map $\nabla:M\to P\otimes_A M$ satisfying
$\nabla(am)=d_0a\otimes m+a\nabla m$.  Its extension is
\[
 d_M(\omega\otimes m)=d\omega\otimes m+
          (-1)^q\omega\wedge\nabla m,
 \qquad\omega\in\Lambda_A^qP.
\]
It is integrable when $d_M^2=0$, equivalently when this square vanishes
on $M$.  The sign is forced by the graded module rule.
For a formal groupoid $G$, a stratification is a unital isomorphism
between the two completed pullbacks of $M$ to $G$ satisfying the cocycle
identity on composable pairs.

\begin{theorem}[Proved: coefficients, with no finiteness condition on the module]\label{cat:L6}
For the corresponding flat formal groupoid and Lie coalgebroid of
Theorem~\ref{cat:L5}, differentiation gives an equivalence between
stratifications on arbitrary $A$-modules and integrable $P$-connections.
It preserves tensor products and its underlying functor on $A$-modules
is the identity.
\end{theorem}
\begin{proof}
We give the square-zero argument, which also explains why $M$ need
not be flat.  Set $A_M=A\oplus M$, with $M^2=0$, and
$P_M=A_M\otimes_A P$; this module is flat over $A_M$.
An integrable connection defines a differential on
$\Lambda_{A_M}P_M=A_M\otimes_A\Lambda_AP$ by extending $d$ and
setting $d(m)=\nabla m$.  The connection identity is exactly the
compatibility with the scalar relations, and its zero curvature is
exactly the additional square-zero equation on $M$.
Thus it is a flat Lie coalgebroid over $A_M$.  The inclusion and
projection between $A$ and $A_M$ are compatible coalgebroid maps.

Integrate this coalgebroid by Theorem~\ref{cat:L5}; write its arrow
algebra as $B_M$.  The maps just described give a retraction
$B\to B_M\to B$.  The source map induces
\[
 B\widehat\otimes_{A,s}A_M\longrightarrow B_M.
\]
This is an isomorphism: on associated graded it is the canonical
isomorphism
$A_M\otimes_A\operatorname{Sym}_AP
 =\operatorname{Sym}_{A_M}P_M$.
Here base change on the associated graded is justified even though
$M$ is not flat, because the finite quotients $B/I^n$ are flat by
Lemma~\ref{cat:L3}; tensoring each successive filtration sequence
therefore remains exact.  Complete filtered comparison then applies.
The kernel of $B_M\to B$ is consequently
$B\widehat\otimes_{A,s}M$, a square-zero ideal.
Restricting the target map $A_M\to B_M$ to $M$ gives the coaction
\[
 M\longrightarrow B\widehat\otimes_{A,s}M.
\]
Its counit and coassociativity are exactly identity and composition
for the integrated groupoid.  Inversion makes the corresponding map
of the two pullbacks an isomorphism, hence a stratification.
Its first-order term is $\nabla$ by the integration equivalence.

Conversely a stratification lets $G$ act on
$\operatorname{Spec}(A\oplus M)$ and forms its action groupoid.
Its arrow algebra, as an algebra for the source map, is
$B\widehat\otimes_A(A\oplus M)$, and its conormal is
$(A\oplus M)\otimes_AP$.  Differentiation produces precisely the
square-zero coalgebroid above and hence an integrable connection.
The two operations are inverse by Theorem~\ref{cat:L5}'s full
faithfulness.  An $A$-linear map $M\to N$ is a map of square-zero
algebras over and under $A$; it is compatible with the coalgebroid
structures exactly when it is horizontal.  The same argument proves
full faithfulness on coefficient morphisms.  Finally, differentiating
a tensor product gives the Leibniz sum of the two connections.  Full
faithfulness identifies its integrated stratification with the tensor
product stratification, proving tensor compatibility.
\end{proof}

\begin{proposition}[Proved: the finite-order and derived formal Poincare lemma]\label{cat:L7}
Filter the relative de Rham complex of the source map by total order,
where an augmentation-ideal factor and a one-form both have order one.
Let $K_n$ be its quotient by total order greater than $n$.
The augmentation $K_n\to A$ is a quasi-isomorphism for every $n$.
The same statement holds for the linearized complex of any stratified
coefficient module, with augmentation to that module.
Consequently
\[
 R\!\varprojlim_n K_n\simeq A
\]
in the derived category of sheaves, and likewise with coefficients.
For modules over a ring the ordinary inverse-limit complex also
computes this derived limit.  For arbitrary sheaves that last assertion
requires a derived-limit-acyclicity hypothesis on the degreewise
towers and is not asserted without it.
\end{proposition}
\begin{proof}
Translation in the groupoid identifies relative one-forms with the
completed pullback of $P$: restriction to the identity and translation
back are inverse by the groupoid identities.  On total order $N$,
Lemma~\ref{cat:L3} therefore identifies the graded complex with
\[
 \left(\operatorname{Sym}_A^{N-q}P\otimes_A\Lambda_A^qP\right)_{0\leq q\leq N},
 \qquad d(u)=du,\quad d(du)=0.
\]
Let $\iota$ be the degree-minus-one derivation on
$\operatorname{Sym}_AP\otimes_A\Lambda_AP$ that sends $du$ to $u$
and kills polynomial generators.  The graded commutator satisfies
$d\iota+\iota d=N$ on total order $N$: both sides are derivations
and agree on the generators.  For $N>0$, division by $N$ gives a
canonical contracting homotopy, defined without a basis or any choice.
Order zero is $A$ in degree zero.
The complex $K_n$ has a finite filtration by these graded pieces.
The long exact sequences of cohomology for its successive short exact
sequences prove $K_n\to A$ a quasi-isomorphism by finite induction.
This argument is valid inside the topos and makes no infinite sequence
of local choices.

A stratification identifies the two pulled-back coefficient modules.
Use that identification to trivialize the coefficient connection along
the source.  The graded homotopy above tensored with $M$ remains a
contracting homotopy; no flatness of $M$ is needed.  The same finite
filtration argument proves the coefficient statement.
The augmentations are compatible with the transition maps.
Applying the derived inverse-limit functor to this levelwise
quasi-isomorphism gives the claimed isomorphism with the derived limit
of the constant system $A$, which is $A$.
For ordinary modules, the degreewise transition maps are surjective.
A countable surjective tower of modules has no higher derived limits:
the standard difference map $\prod M_n\to\prod M_n$ is surjective
by recursively lifting its coordinates, and its kernel is $\lim M_n$.
Thus the ordinary complex computes the derived limit in that case.
This recursive argument is not valid for arbitrary sheaf epimorphisms;
we have expressly avoided using it there.
\end{proof}

\begin{proposition}[Proved: the precise connection with the crystalline starting theorem]\label{cat:L8}
For a flat formal groupoid, let its formal descent site mean the site
of the nilpotent jets of its simplicial nerve, with all simplicial
and jet-transition maps as morphisms, and coverings induced by the
chosen topology on the base.
A crystal on this site means a module cartesian for these transition
maps.  Such crystals are exactly stratified modules and hence, by
Theorem~\ref{cat:L6}, modules with integrable coalgebroid connection.
The finite-order resolution and its derived completion are those of
Proposition~\ref{cat:L7}.

More concretely, if $X\to S$ is a smooth morphism of
$\mathbf Q$-schemes and the groupoid is the formal completion of
$X\times_SX$ along its diagonal, this equivalence is the usual
crystalline starting equivalence: crystals of $\mathcal O$-modules
on the nilpotent infinitesimal site of $X/S$ are equivalent to
$\mathcal O_X$-modules with integrable $S$-connection.
In characteristic zero the nilpotent divided-power crystalline site
has the same objects, because divided powers are uniquely
$x^{[n]}=x^n/n!$.
\end{proposition}
\begin{proof}
A cartesian module on the nerve restricts in degree zero to an
$A$-module $M$.  The two maps from degree one to degree zero identify
its degree-one value with the two pullbacks of $M$; their comparison
is a stratification at each jet order.  The degree-two cartesian
compatibility is its cocycle identity, and the degeneracy is its
unit identity.  Conversely these identities identify the value at
every degree and every jet with the appropriate pullback of $M$;
associativity gives all higher compatibilities.  These two
constructions are mutually inverse on objects and on morphisms.
Theorem~\ref{cat:L6} proves the first equivalence, and
Proposition~\ref{cat:L7} supplies its formal Poincare resolution in
the stated pro/derived sense.

For the assertion about the usual infinitesimal site, an object is a
nilpotent closed immersion $U\hookrightarrow T$ over $S$, with $U$
open in $X$; coverings are open coverings of $T$.
Smoothness gives locally a lift $f:T\to X$ of $U\to X$.
The pullback $f^*M$ is independent of that lift in the required
canonical sense: two lifts $f,g$ agree on $U$, so the pair $(f,g)$
factors through a finite jet of the diagonal formal groupoid, and
its pullback of the stratification gives
$f^*M\simeq g^*M$.  For three lifts the cocycle identity gives the
compatibility of these identifications.  They therefore glue over
open covers of $T$, are independent of the local lifts and are
compatible with all morphisms of infinitesimal thickenings.  We obtain
a crystal.  In the reverse direction evaluate a crystal on $X$ and
on the diagonal jets; its cartesian isomorphisms give the
stratification.  Local lifting shows these operations are inverse,
because every value on $(U,T)$ is recovered as $f^*M$.
The same argument applies to module morphisms and tensor products.

For completeness the divided-power assertion uses the identity
$n!\gamma_n(x)=x^n$, forced by the divided-power axioms.  Since $n!$
is invertible, it forces the displayed formula; that formula satisfies
all axioms and is finite on nilpotent elements.  Hence adding divided
powers supplies no extra datum here.
The margin on 162-5/031 is thus justified by an actual formal descent
site and, for smooth schemes, by the ordinary crystalline site.  The
arbitrary-flat integration theorem itself uses coalgebroids as in
Theorem~\ref{cat:L5}; no unjustified ordinary inverse-limit exactness
in an arbitrary topos is needed.
\end{proof}
